% Motor: XeLaTeX, fijado en .latexmkrc y en la receta del editor.
\documentclass[aspectratio=169,10pt,c]{beamer}

% !TEX root = slides.tex
% Preambulo comun: estilo UTP, colores, fuentes, plantillas de beamer y
% comandos de apoyo didactico. Se carga desde slides.tex tras \documentclass.

% La opción "c" de beamer reparte el espacio libre 1:1.5 (arriba:abajo), por lo
% que el contenido queda por encima del centro. Igualamos ambos rellenos.
\makeatletter
\define@key{beamerframe}{c}[true]{%
  \beamer@frametopskip=0pt plus 1fill\relax%
  \beamer@framebottomskip=0pt plus 1fill\relax%
  \beamer@frametopskipautobreak=0pt plus .5\paperheight\relax%
  \beamer@framebottomskipautobreak=0pt plus .5\paperheight\relax%
}
\makeatother

\usepackage{graphicx}
\usepackage{tikz}
\usepackage{xcolor}
\usepackage{iftex}

\ifPDFTeX
  \usepackage[utf8]{inputenc}
  \usepackage[T1]{fontenc}
  \usepackage{lmodern}          % tipografia escalable para el respaldo con pdflatex
  \renewcommand{\familydefault}{\sfdefault}
\else
  \usepackage{fontspec}
  \IfFontExistsTF{Arial}
    {\setsansfont{Arial}[BoldFont={Arial Bold}]}
    {\setsansfont{Helvetica}[BoldFont={Helvetica Bold}]}
\fi

\usepackage[spanish,es-nodecimaldot]{babel}
\usepackage{microtype}
\usepackage{booktabs}
\usepackage{amssymb}
\usepackage{array}
\usetikzlibrary{arrows.meta,positioning,fit,calc,shapes.geometric}

\graphicspath{{assets/}}

% ---------- Identidad visual UTP ----------
\definecolor{utpRojo}{HTML}{C42B35}
\definecolor{utpNegro}{HTML}{111111}
\definecolor{utpGris}{HTML}{5F5F5F}
\definecolor{utpGrisMedio}{HTML}{9A9A9A}
\definecolor{utpGrisClaro}{HTML}{F1F1F1}
\definecolor{utpNaranja}{HTML}{C1670B}

\usetheme{default}
\setbeamercolor{normal text}{fg=utpNegro,bg=white}
\setbeamercolor{structure}{fg=utpRojo}
\setbeamercolor{frametitle}{fg=utpNegro,bg=white}
\setbeamercolor{title}{fg=utpNegro}
\setbeamercolor{subtitle}{fg=utpNegro}
\setbeamercolor{author}{fg=utpNegro}
\setbeamercolor{date}{fg=utpNegro}
\setbeamercolor{block title}{fg=white,bg=utpNegro}
\setbeamercolor{block body}{fg=utpNegro,bg=utpGrisClaro}
\setbeamercolor{alertblock title}{fg=white,bg=utpRojo}
\setbeamercolor{alertblock body}{fg=utpNegro,bg=utpRojo!7}
\setbeamercolor{exampleblock title}{fg=white,bg=utpGris}
\setbeamercolor{exampleblock body}{fg=utpNegro,bg=utpGris!8}

\setbeamerfont{title}{series=\bfseries,size=\fontsize{19.5}{22.5}\selectfont}
\setbeamerfont{subtitle}{size=\fontsize{11.5}{13.5}\selectfont}
\setbeamerfont{date}{series=\bfseries,size=\fontsize{11.5}{13.5}\selectfont}
\setbeamerfont{author}{size=\fontsize{10.5}{12.5}\selectfont}
\setbeamerfont{frametitle}{series=\bfseries,size=\large}
\setbeamerfont{block title}{series=\bfseries,size=\small}
\setbeamerfont{block body}{size=\normalsize}

\setbeamertemplate{navigation symbols}{}
\setbeamertemplate{headline}{}
\setbeamertemplate{blocks}[default]
\setbeamertemplate{itemize item}{\color{utpRojo}\small\raise.15ex\hbox{$\blacktriangleright$}}
\setbeamertemplate{itemize subitem}{\color{utpGris}\scriptsize$\bullet$}
\setbeamertemplate{enumerate item}{\color{utpRojo}\insertenumlabel.}
\setbeamertemplate{background canvas}{\color{white}\rule{\paperwidth}{\paperheight}}

% Título de diapositiva: negrita negra a la izquierda, logo UTP a la derecha
\setbeamertemplate{frametitle}{%
  \begin{tikzpicture}[remember picture,overlay]
    \node[anchor=north east,inner sep=0pt]
      at ([xshift=-7mm,yshift=-5.5mm]current page.north east)
      {\includegraphics[width=17mm]{utp-logo.png}};
  \end{tikzpicture}%
  \vspace{1.2mm}%
  \begin{minipage}[t]{.80\textwidth}\raggedright\insertframetitle\end{minipage}\par
  \vspace{1.1mm}{\color{utpRojo}\rule{\textwidth}{1.1pt}}
}

% Solo numeración discreta; la plantilla UTP no lleva pie de página
\setbeamertemplate{footline}{%
  \hfill{\color{utpGrisMedio}\scriptsize\insertframenumber}\hspace{6mm}\vspace{2.5mm}%
}

% Portada UTP
\setbeamertemplate{title page}{%
  \begin{tikzpicture}[remember picture,overlay]
    \node[anchor=center,inner sep=0pt] at (current page.center) {%
      \includegraphics[width=\paperwidth,height=\paperheight]{utp-cover-bg.jpeg}%
    };
    \node[anchor=north,align=center,text=utpNegro,inner sep=0pt,
      text width=145mm,font=\usebeamerfont{title}]
      at ([yshift=-25.2mm]current page.north) {\inserttitle};
    \node[anchor=north,align=center,text=utpNegro,inner sep=0pt,
      text width=142mm,font=\usebeamerfont{subtitle}]
      at ([yshift=-38.2mm]current page.north) {\insertsubtitle};
    \node[anchor=north,align=center,text=utpNegro,inner sep=0pt,
      font=\usebeamerfont{date}]
      at ([yshift=-50.4mm]current page.north) {\insertdate};
    \node[anchor=north,align=center,text=utpNegro,inner sep=0pt,
      font=\usebeamerfont{author}]
      at ([yshift=-58.4mm]current page.north) {\insertauthor};
  \end{tikzpicture}%
}

% ---------- Comandos de apoyo didáctico ----------
\newcommand{\concepto}[2]{%
  \begin{block}{#1}\large #2\end{block}}
\newcommand{\pregunta}[1]{%
  \begin{alertblock}{Decide y explica}#1\end{alertblock}}
\newcommand{\analogia}[1]{%
  \begin{exampleblock}{En palabras simples}\small #1\end{exampleblock}}
\newcommand{\enlarealidad}[1]{%
  \begin{block}{Cómo se ve en la realidad}\small #1\end{block}}
\newcommand{\ojo}[1]{%
  \begin{alertblock}{Ojo con esto}\small #1\end{alertblock}}
\newcommand{\apoyo}{{\normalfont\footnotesize\color{utpGris}~· material de apoyo}}
\newcommand{\marca}[1]{\tikz[baseline=(m.base)]{%
  \node[rounded corners=1.5pt,fill=utpGrisClaro,draw=utpGrisMedio,inner xsep=3pt,inner ysep=1.5pt,
        font=\footnotesize] (m) {#1};}}

\tikzset{
  caja/.style={draw=utpNegro, very thick, rounded corners=1.5pt, fill=white,
    align=center, minimum height=0.72cm, inner sep=5pt, font=\small},
  cajaazul/.style={caja, draw=utpGris, fill=white},
  cajaverde/.style={caja, draw=utpNegro, fill=utpGrisClaro},
  cajanaranja/.style={caja, draw=utpRojo!75, fill=utpRojo!7},
  flecha/.style={-{Stealth[length=2.3mm]}, very thick, draw=utpNegro},
  flechavuelta/.style={{Stealth[length=2.3mm]}-{Stealth[length=2.3mm]}, very thick, draw=utpNegro}
}

% ---------- Notación pata de gallo (crow's foot) ----------
% La rotación indica hacia dónde está la entidad: 0 derecha, 180 izquierda,
% 90 arriba, -90 abajo. La pata se abre hacia la entidad del lado "muchos".
\tikzset{
  relacion/.style={thick, draw=utpNegro},
  pics/uno/.style={code={%
    \draw[thick,utpNegro,line cap=round] (0,-1.5mm) -- (0,1.5mm);}},
  pics/muchos/.style={code={%
    \draw[thick,utpNegro,line cap=round,line join=round]
      (2.6mm,1.5mm) -- (0,0) -- (2.6mm,-1.5mm);
    \draw[thick,utpNegro,line cap=round] (0,0) -- (2.6mm,0);}},
}
\newcommand{\simbuno}{\tikz[baseline=-0.6ex]{\draw[relacion](0,0)--(7mm,0);\pic at (4.4mm,0){uno};}}
\newcommand{\simbmuchos}{\tikz[baseline=-0.6ex]{\draw[relacion](0,0)--(7mm,0);\pic at (4.4mm,0){muchos};}}

\usepackage[normalem]{ulem}

\title{Vistas, DBML y DDL en SQLite}
\subtitle{Consultas guardadas, diagramas como código,\\ el DDL en SQLite y los archivos \texttt{.sql}}
\date{Semana 08 - Sesión 16}
\author{Rel Guzman}

% Espacio en blanco para completar
\newcommand{\bl}[1][1.9cm]{\underline{\hspace{#1}}}
% Nombre de tabla, siempre en naranja
\newcommand{\tbl}[1]{\textcolor{utpNaranja}{\texttt{#1}}}
% Palabra reservada de SQL o de DBML dentro de un bloque de código
\newcommand{\kw}[1]{\textcolor{utpRojo}{\textbf{#1}}}
% Parte nueva de un bloque de codigo que crece paso a paso
\newcommand{\nuevo}[1]{{%
  \setlength{\fboxsep}{1.2pt}%
  \begingroup
  \renewcommand{\kw}[1]{\textcolor{utpNegro}{\textbf{##1}}}%
  \colorbox{yellow!45}{\textcolor{utpNegro}{#1}}%
  \endgroup}}
% Comentario dentro de un bloque de código SQL
\newcommand{\com}[1]{\textcolor{utpGris}{-- #1}}
% Comentario dentro de un bloque de código DBML
\newcommand{\comd}[1]{\textcolor{utpGris}{// #1}}
% Marcas de clave primaria y clave foránea
\newcommand{\pk}{{\scriptsize\color{utpRojo}\textbf{PK}}}
\newcommand{\fk}{{\scriptsize\color{utpNaranja}\textbf{FK}}}
% Insignia de práctica
\newcommand{\ins}[1]{\tikz[baseline=-0.5ex]{\node[rounded corners=2pt,fill=utpRojo,text=white,
  font=\bfseries\scriptsize,inner xsep=5pt,inner ysep=2.5pt] {#1};}}
% Etiqueta oscura: nombre del gestor o de la forma de escribir algo
\newcommand{\etiqueta}[1]{\tikz[baseline=-0.5ex]{\node[rounded corners=2pt,fill=utpNegro,text=white,
  font=\bfseries\scriptsize,inner xsep=6pt,inner ysep=3pt] {#1};}}
% Celda con un dato que rompe una regla
\newcommand{\mal}[1]{{\setlength{\fboxsep}{1.2pt}\colorbox{utpRojo!18}{#1}}}
% Valor desconocido
\newcommand{\nulo}{\textcolor{utpGrisMedio}{\textit{NULL}}}
% Se puede / no se puede
\newcommand{\si}{\textbf{sí}}
\newcommand{\no}{\textcolor{utpRojo}{\textbf{no}}}
% Bloque de código sobre fondo gris
\newsavebox{\sqlbuf}
\newenvironment{sqlcode}[1][.86\textwidth]
  {\begin{center}\begin{lrbox}{\sqlbuf}\begin{minipage}{#1}\ttfamily\footnotesize\vspace{6pt}\raggedright}
  {\vspace{6pt}\end{minipage}\end{lrbox}\colorbox{utpGrisClaro}{\usebox{\sqlbuf}}\end{center}}
% Sangría dentro de un bloque de código
\newcommand{\ind}{\hspace*{1.6em}}
% Una forma de escribir algo: etiqueta a la izquierda, código a la derecha
\newcommand{\forma}[3]{%
  \begin{columns}[c,onlytextwidth]
  \column{.30\textwidth}\etiqueta{#1}\\[3pt]{\scriptsize\color{utpGris} #2}
  \column{.68\textwidth}\begin{sqlcode}[\textwidth]\scriptsize #3\end{sqlcode}
  \end{columns}\vspace{-0.35cm}}
% Anatomía de una instrucción: cada pieza del código con su pregunta al lado
\tikzset{
  pieza/.style={fill=utpGrisClaro,font=\ttfamily\footnotesize,anchor=west,inner sep=5pt,
    minimum height=0.62cm},
  para/.style={font=\small,anchor=west,text=utpNegro}
}
\newcommand{\piezax}{7.6}
\newcommand{\pieza}[3]{%
  \node[pieza] at (0,#1) {#2};
  \node[para] at (\piezax,#1) {\textcolor{utpRojo}{$\blacktriangleleft$}\ #3};}
% Diapositiva de sección: una palabra bien grande (\seccion[0]{...} la deja sin línea roja)
\newcommand{\seccion}[2][1]{%
\begin{frame}[plain]
  \begin{tikzpicture}[remember picture,overlay]
    \node[anchor=north east,inner sep=0pt]
      at ([xshift=-10mm,yshift=-8mm]current page.north east)
      {\includegraphics[width=43mm]{utp-logo.png}};
    \node[anchor=center,inner sep=0pt] at (current page.center)
      {\resizebox{100mm}{!}{\bfseries\color{utpNegro}#2}};
    \ifnum#1=1
    \node[anchor=north,inner sep=0pt] at ([yshift=-11mm]current page.center)
      {\color{utpRojo}\rule{100mm}{1.6pt}};
    \fi
  \end{tikzpicture}
\end{frame}}
% Encabezado y filas de la tabla de ejemplo
\newcommand{\cabproducto}{%
  \normalfont\bfseries id & \normalfont\bfseries nombre & \normalfont\bfseries categoria
  & \normalfont\bfseries precio & \normalfont\bfseries stock \\}
\newcommand{\filasproducto}{%
  1 & Arroz & abarrotes & 4.50 & 30 \\
  2 & Gaseosa & bebidas & 3.00 & 12 \\
  3 & Agua & bebidas & 1.50 & 0 \\
  4 & Galletas & snacks & 1.20 & 50 \\
  5 & Leche & bebidas & 3.80 & 20 \\
  6 & Chifles & \nulo & 2.00 & 8 \\}
% Tabla de ejemplo completa, lista para poner al costado
\newcommand{\tablaproducto}{%
  \centering\tbl{producto}\\[4pt]
  \scriptsize\ttfamily
  \begin{tabular}{lllrr}
  \toprule
  \cabproducto
  \midrule
  \filasproducto
  \bottomrule
  \end{tabular}}
% Diagrama de la tienda: un cliente hace muchos pedidos (horizontal y vertical)
\newcommand{\diagtienda}{%
  \begin{tikzpicture}
    \node[cajaverde,text width=1.9cm,font=\scriptsize,align=left] (cli)
      {\textbf{cliente}\\ \texttt{id} \pk\\ \texttt{nombre}\\ \texttt{correo}};
    \node[cajanaranja,text width=2.2cm,font=\scriptsize,align=left,right=13mm of cli] (ped)
      {\textbf{pedido}\\ \texttt{id} \pk\\ \texttt{fecha}\\ \texttt{cliente\_id} \fk};
    \draw[relacion] (cli)--(ped);
    \pic at ($(cli.east)!2.6mm!(ped.west)$) {uno};
    \pic at ($(ped.west)!2.6mm!(cli.east)$) {muchos};
  \end{tikzpicture}}
\newcommand{\diagtiendav}{%
  \begin{tikzpicture}
    \node[cajaverde,text width=2.2cm,font=\scriptsize,align=left] (cli)
      {\textbf{cliente}\\ \texttt{id} \pk\\ \texttt{nombre}\\ \texttt{correo}};
    \node[cajanaranja,text width=2.2cm,font=\scriptsize,align=left,below=11mm of cli] (ped)
      {\textbf{pedido}\\ \texttt{id} \pk\\ \texttt{fecha}\\ \texttt{cliente\_id} \fk};
    \draw[relacion] (cli)--(ped);
    \pic[rotate=90] at ($(cli.south)!2.6mm!(ped.north)$) {uno};
    \pic[rotate=-90] at ($(ped.north)!2.6mm!(cli.south)$) {muchos};
  \end{tikzpicture}}

\begin{document}

% 1 ---------------------------------------------------------------------
{\setbeamertemplate{footline}{}\maketitle}

% 2 ---------------------------------------------------------------------
\begin{frame}[c]
\begin{block}{Logro de la sesión}
Al finalizar, podrás \textbf{crear y borrar una vista}, \textbf{escribir el diseño} de una base de datos en \textbf{DBML} para ver su diagrama, reconocer \textbf{qué cambia en SQLite} al crear, modificar y borrar tablas, y explicar \textbf{qué es un archivo \texttt{.sql}}.
\end{block}
\end{frame}
% ---------------------------------------------------------------------
\begin{frame}{Lo de hoy: cinco paradas}
\begin{center}
\begin{tikzpicture}
  \tikzset{parada/.style={cajaverde,font=\footnotesize\bfseries,text width=1.95cm,minimum height=1.15cm}}
  \node[parada] (a) {1. Vistas};
  \node[parada,right=4.5mm of a] (b) {2. Diagramas\\como código};
  \node[parada,right=4.5mm of b] (c) {3. DBML};
  \node[parada,right=4.5mm of c] (d) {4. DDL\\en SQLite};
  \node[parada,right=4.5mm of d] (e) {5. Archivos\\\texttt{.sql}};
  \draw[flecha] (a)--(b); \draw[flecha] (b)--(c); \draw[flecha] (c)--(d); \draw[flecha] (d)--(e);
\end{tikzpicture}
\end{center}
\vspace{0.5cm}
\begin{itemize}\setlength\itemsep{8pt}
  \item Ya sabes \textbf{crear tablas} y \textbf{consultarlas}. Hoy no hay instrucciones difíciles.
  \item Hoy vemos \textbf{ayudas}: guardar una consulta, dibujar el diseño escribiendo y guardar todo en un archivo.
\end{itemize}
\vspace{0.4cm}
\ojo{Todo es \textbf{introductorio}. La meta es entender \textbf{la idea} de cada cosa.}
\end{frame}
% =======================================================================
\seccion{Vistas}
% ---------------------------------------------------------------------
\begin{frame}{La misma consulta, una y otra vez}
\begin{columns}[c,onlytextwidth]
\column{.46\textwidth}
\tablaproducto
\column{.50\textwidth}
\small El encargado de bebidas pide esto \textbf{todos los días}:
\begin{sqlcode}[\textwidth]
\scriptsize
\kw{SELECT} nombre, precio\\
\kw{FROM} producto\\
\kw{WHERE} categoria = 'bebidas';
\end{sqlcode}
\vspace{-0.2cm}
\small Escribirla cada vez \textbf{cansa}, y es fácil \textbf{equivocarse}.
\end{columns}
\vspace{0.3cm}
\concepto{La idea}{\small Le ponemos \textbf{un nombre} a esa consulta y el gestor \textbf{la guarda}. Eso es una \textbf{vista}.}
\vspace{0.1cm}
\analogia{Es como \textbf{guardar un contacto}: en vez de marcar el número completo cada vez, buscas «Mamá».}
\end{frame}
% ---------------------------------------------------------------------
\begin{frame}{\texttt{CREATE VIEW}: una consulta con nombre}
\begin{center}
\begin{tikzpicture}[y=0.78cm]
  \pieza{2}{\nuevo{\kw{CREATE VIEW}} bebidas \nuevo{\kw{AS}}}{el \textbf{nombre} que le pones}
  \pieza{1}{\kw{SELECT} nombre, precio \kw{FROM} producto}{la \textbf{consulta} que se guarda}
  \pieza{0}{\kw{WHERE} categoria = 'bebidas';}{con su condición y todo}
\end{tikzpicture}
\end{center}
\vspace{0.25cm}
\begin{columns}[c,onlytextwidth]
\column{.56\textwidth}
Desde ahora se usa \textbf{como si fuera una tabla}:
\begin{sqlcode}[\textwidth]
\kw{SELECT} * \kw{FROM} \nuevo{bebidas};
\end{sqlcode}
\column{.38\textwidth}
\centering\footnotesize\ttfamily
\begin{tabular}{lr}
\toprule
\normalfont\bfseries nombre & \normalfont\bfseries precio \\
\midrule
Gaseosa & 3.00 \\
Agua & 1.50 \\
Leche & 3.80 \\
\bottomrule
\end{tabular}
\end{columns}
\vspace{0.35cm}
\small El \texttt{CREATE VIEW} se escribe \textbf{una sola vez}. Es igual en MySQL, SQL Server y SQLite.
\end{frame}
% ---------------------------------------------------------------------
\begin{frame}{Una vista no guarda datos: es una ventana}
\begin{center}
\begin{tikzpicture}
  \node[cajaverde,text width=3.5cm] (t) {\textbf{Tabla} \tbl{producto}\\[-1pt]{\scriptsize guarda las \textbf{filas}}};
  \node[caja,dashed,text width=3.5cm,right=30mm of t] (v) {\textbf{Vista} \texttt{bebidas}\\[-1pt]{\scriptsize guarda solo la \textbf{consulta}}};
  \draw[flecha] (v)--(t) node[midway,above,font=\scriptsize] {cada vez que la usas,} node[midway,below,font=\scriptsize] {va a leer la tabla};
\end{tikzpicture}
\end{center}
\vspace{0.15cm}
\begin{columns}[c,onlytextwidth]
\column{.56\textwidth}
\small Llega un producto nuevo a la \textbf{tabla}:
\begin{sqlcode}[\textwidth]
\scriptsize
\kw{INSERT INTO} producto\\
\kw{VALUES} (7, 'Chicha', 'bebidas', 2.50, 15);
\end{sqlcode}
\vspace{-0.2cm}
\small La Chicha aparece \textbf{sola} en la vista. Nadie tocó \texttt{bebidas}.
\column{.38\textwidth}
\centering\scriptsize\ttfamily
\begin{tabular}{lr}
\toprule
\normalfont\bfseries nombre & \normalfont\bfseries precio \\
\midrule
Gaseosa & 3.00 \\
Agua & 1.50 \\
Leche & 3.80 \\
\nuevo{Chicha} & \nuevo{2.50} \\
\bottomrule
\end{tabular}
\end{columns}
\vspace{0.3cm}
\analogia{Una \textbf{ventana} no guarda el paisaje: te muestra lo que hay afuera \textbf{en este momento}.}
\end{frame}
% ---------------------------------------------------------------------
\begin{frame}{\texttt{DROP VIEW}: borrar la vista, no los datos}
\begin{columns}[T,onlytextwidth]
\column{.48\textwidth}
\begin{sqlcode}[\textwidth]
\nuevo{\kw{DROP VIEW}} bebidas;
\end{sqlcode}
\vspace{-0.2cm}
Se borra \textbf{la consulta guardada}.\\[3pt]
La tabla \tbl{producto} sigue con \textbf{todas sus filas}.
\column{.48\textwidth}
\begin{sqlcode}[\textwidth]
\kw{DROP TABLE} producto;
\end{sqlcode}
\vspace{-0.2cm}
Se borra \textbf{la tabla con sus datos}.\\[3pt]
Esto sí es grave.
\end{columns}
\vspace{0.5cm}
\analogia{Borrar un \textbf{contacto} del celular no desaparece a la persona. Borrar una vista no desaparece los datos.}
\vspace{0.2cm}
\small ¿Y si quiero \textbf{cambiar} una vista? Lo más simple: la \textbf{borras} y la \textbf{creas de nuevo}.
\end{frame}
% ---------------------------------------------------------------------
\begin{frame}{¿Cuándo es útil una vista?}
\begin{columns}[T,onlytextwidth]
\column{.31\textwidth}
\centering\etiqueta{Para no repetir}\\[6pt]
\raggedright\small Una consulta que se usa seguido se escribe \textbf{una vez}.
\column{.31\textwidth}
\centering\etiqueta{Para simplificar}\\[6pt]
\raggedright\small Detrás puede haber algo largo; quien la usa ve algo \textbf{sencillo}.
\column{.31\textwidth}
\centering\etiqueta{Para ocultar}\\[6pt]
\raggedright\small Muestra solo \textbf{algunas columnas}: el precio sí, el costo no.
\end{columns}
\vspace{0.6cm}
\begin{alertblock}{¿Y cuándo no?}
\small Si la consulta la vas a usar \textbf{una sola vez}, no hace falta una vista.
\end{alertblock}
\vspace{0.2cm}
\enlarealidad{En un colegio, la secretaría usa una vista con el \textbf{nombre y el teléfono} de cada alumno. Las notas y los pagos no aparecen ahí.}
\end{frame}
% ---------------------------------------------------------------------
\begin{frame}{\ins{Práctica 1} Vistas sobre la bodega}
\begin{columns}[c,onlytextwidth]
\column{.44\textwidth}
\tablaproducto
\column{.52\textwidth}
\small Se crea esta vista:
\begin{sqlcode}[\textwidth]
\scriptsize
\kw{CREATE VIEW} baratos \kw{AS}\\
\kw{SELECT} nombre, precio\\
\kw{FROM} producto\\
\kw{WHERE} precio < 3;
\end{sqlcode}
\end{columns}
\vspace{0.2cm}
\begin{center}
\small
\begin{tabular}{>{\raggedright\arraybackslash}p{8.6cm}>{\raggedright\arraybackslash}p{3.4cm}}
\toprule
\bfseries Pregunta & \bfseries Respuesta \\
\midrule
¿Qué nombres muestra \texttt{SELECT * FROM baratos;}? & \bl[3.2cm] \\[9pt]
La Gaseosa baja a 2.50. ¿Aparece ahora en \texttt{baratos}? & \bl[3.2cm] \\[9pt]
Se ejecuta \texttt{DROP VIEW baratos;}. ¿Cuántas filas tiene \tbl{producto}? & \bl[3.2cm] \\
\bottomrule
\end{tabular}
\end{center}
\vspace{0.2cm}
\small \textbf{Para escribir:} crea la vista \texttt{agotados}, con los productos cuyo stock es 0.
\end{frame}
% =======================================================================
\seccion{Diagramas como código}
% ---------------------------------------------------------------------
\begin{frame}{Dos maneras de hacer un diagrama}
\begin{columns}[T,onlytextwidth]
\column{.47\textwidth}
\centering\etiqueta{Con el mouse}\\[8pt]
\raggedright
\begin{itemize}\setlength\itemsep{6pt}
  \item Arrastras cajas y las unes con líneas, \textbf{a mano}.
  \item Si agregas una tabla, toca \textbf{reacomodar} todo.
\end{itemize}
\column{.49\textwidth}
\centering\etiqueta{Con código}\\[8pt]
\raggedright
\begin{itemize}\setlength\itemsep{6pt}
  \item \textbf{Escribes} qué tablas hay y cómo se conectan.
  \item Un programa \textbf{dibuja el diagrama} por ti.
\end{itemize}
\end{columns}
\vspace{0.35cm}
\begin{center}
\begin{tikzpicture}
  \node[cajaverde,font=\footnotesize\ttfamily,align=left,text width=3.4cm] (txt) {cliente\\pedido\\pedido > cliente};
  \node[caja,font=\footnotesize,right=38mm of txt,text width=1.7cm] (c) {cliente};
  \node[caja,font=\footnotesize,right=10mm of c,text width=1.7cm] (p) {pedido};
  \draw[relacion] (c)--(p);
  \draw[flecha] (txt)--(c) node[midway,above,font=\scriptsize] {un programa lo dibuja};
  \node[font=\scriptsize,text=utpGris,below=1.5mm of txt] {tú escribes esto};
  \node[font=\scriptsize,text=utpGris] at ($(c.south)!.5!(p.south)+(0,-3.2mm)$) {y obtienes esto};
\end{tikzpicture}
\end{center}
\vspace{0.1cm}
\analogia{Es la diferencia entre la \textbf{foto de un plato} y su \textbf{receta}. En la receta, cambiar un ingrediente es cambiar \textbf{una línea}.}
\end{frame}
% ---------------------------------------------------------------------
\begin{frame}{¿Por qué escribir los diagramas?}
\begin{columns}[c,onlytextwidth]
\column{.54\textwidth}
\begin{itemize}\setlength\itemsep{9pt}
  \item \textbf{Es rápido de cambiar:} corriges una línea y el dibujo se rehace solo.
  \item \textbf{Es texto:} se copia, se pega en un chat y se guarda junto al proyecto.
  \item \textbf{Lo puede escribir una IA:} le cuentas el caso, te da el texto y tú lo revisas.
\end{itemize}
\column{.42\textwidth}
\centering\small
\renewcommand{\arraystretch}{1.35}
\begin{tabular}{>{\bfseries}ll}
\toprule
Lenguaje & \bfseries Sirve para \\
\midrule
Mermaid & diagramas de flujo \\
PlantUML & diagramas UML \\
\nuevo{DBML} & \textbf{bases de datos} \\
\bottomrule
\end{tabular}\\[6pt]
\scriptsize\color{utpGris}hay varios lenguajes para dibujar con texto
\end{columns}
\end{frame}
% =======================================================================
\seccion[0]{DBML}
% ---------------------------------------------------------------------
\begin{frame}{DBML: describir una base de datos con texto}
\begin{columns}[c,onlytextwidth]
\column{.50\textwidth}
\begin{sqlcode}[\textwidth]
\scriptsize
\kw{Table} cliente \{\\
\ind id int [pk]\\
\ind nombre varchar(80)\\
\ind correo varchar(120)\\
\}\\[5pt]
\kw{Table} pedido \{\\
\ind id int [pk]\\
\ind fecha date\\
\ind cliente\_id int\\
\}\\[5pt]
\kw{Ref}: pedido.cliente\_id > cliente.id
\end{sqlcode}
\column{.46\textwidth}
\centering\diagtiendav\\[4pt]
{\scriptsize\color{utpGris}el diagrama que sale de ese texto}
\end{columns}
\vspace{0.15cm}
\small \textbf{DBML} es un lenguaje pequeño y gratuito, hecho solo para \textbf{describir tablas y sus relaciones}. Se lee casi sin haberlo estudiado.
\end{frame}
% ---------------------------------------------------------------------
\begin{frame}{¿Por qué DBML, si ya sabemos SQL?}
\begin{columns}[T,onlytextwidth]
\column{.48\textwidth}
\centering\etiqueta{SQL}
\begin{sqlcode}[\textwidth]
\scriptsize
\kw{CREATE TABLE} pedido (\\
\ind id INT \kw{PRIMARY KEY},\\
\ind fecha DATE \kw{NOT NULL},\\
\ind cliente\_id INT,\\
\ind \kw{FOREIGN KEY} (cliente\_id)\\
\ind\ind \kw{REFERENCES} cliente(id)\\
);
\end{sqlcode}
\column{.48\textwidth}
\centering\etiqueta{DBML}
\begin{sqlcode}[\textwidth]
\scriptsize
\kw{Table} pedido \{\\
\ind id int [pk]\\
\ind fecha date [not null]\\
\ind cliente\_id int\\
\}\\[5pt]
\kw{Ref}: pedido.cliente\_id > cliente.id
\end{sqlcode}
\end{columns}
\vspace{0.1cm}
\begin{itemize}\setlength\itemsep{6pt}
  \item \textbf{Es más corto:} menos palabras y sin comas que se olvidan.
  \item \textbf{Es uno solo:} el SQL cambia de un gestor a otro; el DBML no.
  \item \textbf{Se vuelve diagrama al instante:} sirve para conversar el diseño \textbf{antes} de crear nada.
\end{itemize}
\end{frame}
% ---------------------------------------------------------------------
\begin{frame}{DBML no reemplaza a SQL: lo ayuda}
\begin{center}
\begin{tikzpicture}
  \node[cajaverde,text width=2.2cm] (d) {\textbf{DBML}\\[-1pt]{\scriptsize el diseño, en texto}};
  \node[caja,text width=2.3cm,above right=-2mm and 16mm of d] (g) {\textbf{Diagrama}\\[-1pt]{\scriptsize para conversar}};
  \node[caja,text width=2.3cm,below right=-2mm and 16mm of d] (s) {\textbf{SQL}\\[-1pt]{\scriptsize \texttt{CREATE TABLE}\ldots}};
  \node[cajanaranja,text width=2.7cm,right=16mm of s] (b) {\textbf{Base de datos}\\[-1pt]{\scriptsize tablas y datos reales}};
  \draw[flecha] (d.east)--(g.west);
  \draw[flecha] (d.east)--(s.west);
  \draw[flecha] (s)--(b) node[midway,above,font=\tiny] {se ejecuta};
\end{tikzpicture}
\end{center}
\vspace{0.15cm}
\begin{columns}[T,onlytextwidth]
\column{.48\textwidth}
\etiqueta{DBML sí}\\[5pt]
\raggedright\small
\begin{itemize}\setlength\itemsep{4pt}
  \item Describe tablas y relaciones.
  \item Genera el diagrama.
  \item Genera el \texttt{CREATE TABLE}.
\end{itemize}
\column{.48\textwidth}
\etiqueta{DBML no}\\[5pt]
\raggedright\small
\begin{itemize}\setlength\itemsep{4pt}
  \item \textbf{No crea} la base de datos.
  \item \textbf{No guarda} datos.
  \item \textbf{No consulta}: no hay \texttt{SELECT}.
\end{itemize}
\end{columns}
\vspace{0.3cm}
\analogia{DBML es el \textbf{plano} de la casa; SQL es \textbf{la obra}. El plano ayuda a pensar, pero nadie vive en un plano.}
\end{frame}
% ---------------------------------------------------------------------
\begin{frame}{Una tabla en DBML, pieza por pieza}
\begin{center}
\begin{tikzpicture}[y=0.82cm]
  \renewcommand{\piezax}{5.7}
  \pieza{4}{\nuevo{\kw{Table}} cliente \nuevo{\{}}{\texttt{Table}, el \textbf{nombre} y una llave que abre}
  \pieza{3}{\ \ id int [pk]}{una columna por línea: \textbf{nombre} y \textbf{tipo}}
  \pieza{2}{\ \ nombre varchar(80) \nuevo{[not null]}}{entre \textbf{corchetes}, las reglas}
  \pieza{1}{\ \ correo varchar(120) [unique]}{\textbf{sin coma} al final de la línea}
  \pieza{0}{\nuevo{\}}}{la llave que cierra}
\end{tikzpicture}
\end{center}
\vspace{0.4cm}
El orden es el mismo de \texttt{CREATE TABLE}: \textbf{nombre de la columna, tipo y reglas}.
\end{frame}
% ---------------------------------------------------------------------
\begin{frame}{Las reglas entre corchetes}
\begin{center}
\renewcommand{\arraystretch}{1.4}
\begin{tabular}{>{\leavevmode\color{utpRojo}\ttfamily\bfseries}p{2.6cm}>{\ttfamily\small}p{3.6cm}>{\raggedright\arraybackslash}p{5.2cm}}
\toprule
\normalfont\color{utpNegro}\bfseries En DBML & \normalfont\bfseries En SQL & \bfseries Qué significa \\
\midrule
pk & PRIMARY KEY & identifica a cada fila \\
increment & AUTO\_INCREMENT & el gestor numera solo \\
not null & NOT NULL & no se puede dejar vacío \\
unique & UNIQUE & no se puede repetir \\
\bottomrule
\end{tabular}
\end{center}
\vspace{0.2cm}
\begin{sqlcode}[.60\textwidth]
id int [\nuevo{pk}, \nuevo{increment}]
\end{sqlcode}
\vspace{-0.1cm}
Son las \textbf{mismas reglas} que ya conoces, más cortas. Si hay varias, se separan con \textbf{comas}.
\end{frame}
% ---------------------------------------------------------------------
\begin{frame}{Las relaciones: \texttt{Ref}}
\begin{sqlcode}[.70\textwidth]
\nuevo{\kw{Ref}:} pedido.cliente\_id \nuevo{>} cliente.id
\end{sqlcode}
\vspace{-0.1cm}
\begin{columns}[c,onlytextwidth]
\column{.50\textwidth}
\small
\begin{itemize}\setlength\itemsep{6pt}
  \item Se lee: «\texttt{cliente\_id} de \tbl{pedido} \textbf{apunta} al \texttt{id} de \tbl{cliente}».
  \item Es la \textbf{clave foránea}, en una sola línea.
\end{itemize}
\column{.47\textwidth}
\centering\diagtienda\\[5pt]
{\scriptsize\color{utpGris}un cliente, muchos pedidos}
\end{columns}
\vspace{0.25cm}
\begin{center}
\renewcommand{\arraystretch}{1.25}
\begin{tabular}{>{\leavevmode\color{utpRojo}\ttfamily\bfseries}c>{\raggedright\arraybackslash}p{3.0cm}>{\ttfamily\small}p{6.6cm}}
\toprule
\normalfont\color{utpNegro}\bfseries Símbolo & \bfseries Relación & \normalfont\bfseries Ejemplo \\
\midrule
> & muchos a uno & pedido.cliente\_id > cliente.id \\
- & uno a uno & usuario.id - perfil.usuario\_id \\
\bottomrule
\end{tabular}
\end{center}
\vspace{0.1cm}
\small \textbf{Truco:} el \texttt{>} se abre hacia el lado «\textbf{muchos}», igual que la pata de gallo.
\end{frame}
% ---------------------------------------------------------------------
\begin{frame}{\ins{Práctica 2} Lee este DBML}
\begin{columns}[c,onlytextwidth]
\column{.50\textwidth}
\begin{sqlcode}[\textwidth]
\scriptsize
\kw{Table} categoria \{\\
\ind id int [pk, increment]\\
\ind nombre varchar(80) [unique]\\
\}\\[5pt]
\kw{Table} producto \{\\
\ind id int [pk, increment]\\
\ind nombre varchar(80) [not null]\\
\ind precio decimal(10,2)\\
\ind categoria\_id int\\
\}\\[5pt]
\kw{Ref}: producto.categoria\_id > categoria.id
\end{sqlcode}
\column{.46\textwidth}
\small
\begin{tabular}{>{\raggedright\arraybackslash}p{4.1cm}l}
\toprule
\bfseries Pregunta & \bfseries Resp. \\
\midrule
¿Cuántas tablas hay? & \bl[1.4cm] \\[10pt]
¿Qué columna es clave foránea? & \bl[1.4cm] \\[10pt]
¿Pueden repetirse los nombres de categoría? & \bl[1.4cm] \\[10pt]
¿Cuál tabla está del lado «muchos»? & \bl[1.4cm] \\
\bottomrule
\end{tabular}
\end{columns}
\vspace{0.3cm}
\small \textbf{Dibuja} en tu cuaderno el diagrama: dos cajas y la línea que las une.
\end{frame}
% =======================================================================
\seccion{Herramientas}
% ---------------------------------------------------------------------
\begin{frame}{Herramientas que usan DBML}
\begin{center}
\renewcommand{\arraystretch}{1.5}
\begin{tabular}{>{\raggedright\bfseries}p{3.6cm}>{\raggedright\arraybackslash}p{8.6cm}}
\toprule
Herramienta & \bfseries Para qué sirve \\
\midrule
\nuevo{dbdiagram.io} & Escribes DBML y \textbf{ves el diagrama} al lado. Es la que usaremos. \\
dbdocs.io & Convierte el DBML en una \textbf{página web} que explica la base de datos. \\
Extensiones de VS Code & Muestran el diagrama \textbf{dentro del editor}. \\
\bottomrule
\end{tabular}
\end{center}
\vspace{0.4cm}
DBML es \textbf{abierto}: cualquier herramienta puede leerlo. Tu diseño no queda atrapado en un solo programa.
\end{frame}
% ---------------------------------------------------------------------
\begin{frame}{dbdiagram.io en tres pasos}
\begin{center}
\begin{tikzpicture}
  \node[cajaverde,text width=3.5cm,minimum height=1.5cm] (a) {\textbf{1. Escribe}\\[-1pt]{\scriptsize el DBML, a la izquierda}};
  \node[cajaverde,text width=3.5cm,minimum height=1.5cm,right=9mm of a] (b) {\textbf{2. Mira}\\[-1pt]{\scriptsize el diagrama, a la derecha}};
  \node[cajaverde,text width=3.5cm,minimum height=1.5cm,right=9mm of b] (c) {\textbf{3. Exporta}\\[-1pt]{\scriptsize a SQL, PDF o imagen}};
  \draw[flecha] (a)--(b); \draw[flecha] (b)--(c);
\end{tikzpicture}
\end{center}
\vspace{0.4cm}
\begin{itemize}\setlength\itemsep{7pt}
  \item Se usa desde el navegador: \textbf{no se instala nada}.
  \item Al exportar a SQL eliges el gestor, por ejemplo \textbf{MySQL}.
\end{itemize}
\vspace{0.4cm}
\ojo{En la lista \textbf{no aparece SQLite}: el SQL que sale hay que \textbf{ajustarlo un poco}. Eso vemos en la siguiente parte.}
\end{frame}
% ---------------------------------------------------------------------
\begin{frame}{\ins{Práctica 3} Del enunciado al diagrama}
Un colegio registra a sus \textbf{alumnos} y las \textbf{matrículas} de cada uno.
\vspace{0.3cm}
\begin{columns}[T,onlytextwidth]
\column{.48\textwidth}
\centering\etiqueta{alumno}\\[6pt]
\raggedright\small
\begin{itemize}\setlength\itemsep{3pt}
  \item \texttt{id}: clave primaria
  \item \texttt{nombre}: obligatorio
  \item \texttt{correo}: no se repite
\end{itemize}
\column{.48\textwidth}
\centering\etiqueta{matricula}\\[6pt]
\raggedright\small
\begin{itemize}\setlength\itemsep{3pt}
  \item \texttt{id}: clave primaria
  \item \texttt{fecha}: obligatoria
  \item \texttt{alumno\_id}: apunta al alumno
\end{itemize}
\end{columns}
\vspace{0.4cm}
\begin{enumerate}\setlength\itemsep{6pt}
  \item Escribe en DBML las \textbf{dos tablas} y la \textbf{relación}.
  \item Pégalo en \texttt{dbdiagram.io} y mira el diagrama.
  \item Exporta a \textbf{MySQL} y guarda ese SQL: lo usarás en la práctica 4.
\end{enumerate}
\end{frame}
% =======================================================================
\seccion{DDL en SQLite}
% ---------------------------------------------------------------------
\begin{frame}{¿Por qué SQLite?}
\begin{center}
\begin{tikzpicture}
  \node[caja,text width=2.3cm] (p1) {Tu programa};
  \node[cajaverde,text width=2.7cm,right=9mm of p1] (s1) {\textbf{Servidor}\\[-1pt]{\scriptsize MySQL, SQL Server}};
  \node[caja,text width=1.9cm,right=9mm of s1] (d1) {Datos};
  \draw[flechavuelta] (p1)--(s1); \draw[flechavuelta] (s1)--(d1);
  \node[caja,text width=2.3cm,below=6mm of p1] (p2) {Tu programa};
  \node[cajanaranja,text width=2.7cm,right=9mm of p2] (d2) {\textbf{Un archivo}\\[-1pt]{\scriptsize \texttt{tienda.db}}};
  \draw[flechavuelta] (p2)--(d2);
  \node[font=\scriptsize,text=utpGris,anchor=west] at ([xshift=4mm]d1.east) {hay que instalar y encender};
  \node[font=\scriptsize,text=utpGris,anchor=west] at ([xshift=4mm]d2.east) {\textbf{SQLite}: sin servidor};
\end{tikzpicture}
\end{center}
\vspace{0.2cm}
\begin{itemize}\setlength\itemsep{6pt}
  \item Toda la base de datos es \textbf{un solo archivo}: se copia y se envía como una foto.
  \item \textbf{No se instala} nada ni se crean usuarios.
  \item Está \textbf{en todas partes}: en tu celular, en tu navegador y dentro de Python.
\end{itemize}
\vspace{0.25cm}
\ojo{Por ser tan pequeño, SQLite \textbf{simplifica}. Las instrucciones ya las conoces: solo veremos \textbf{qué cambia}.}
\end{frame}
% ---------------------------------------------------------------------
\begin{frame}{\texttt{CREATE TABLE} en SQLite}
\begin{columns}[T,onlytextwidth]
\column{.48\textwidth}
\centering\etiqueta{MySQL}
\begin{sqlcode}[\textwidth]
\scriptsize
\kw{CREATE TABLE} cliente (\\
\ind id INT \kw{PRIMARY KEY} \kw{AUTO\_INCREMENT},\\
\ind nombre VARCHAR(80) \kw{NOT NULL},\\
\ind registro DATE\\
);
\end{sqlcode}
\column{.48\textwidth}
\centering\etiqueta{SQLite}
\begin{sqlcode}[\textwidth]
\scriptsize
\kw{CREATE TABLE} cliente (\\
\ind id \nuevo{INTEGER} \kw{PRIMARY KEY} \nuevo{\kw{AUTOINCREMENT}},\\
\ind nombre \nuevo{TEXT} \kw{NOT NULL},\\
\ind registro \nuevo{TEXT}\\
);
\end{sqlcode}
\end{columns}
\vspace{0.1cm}
\begin{itemize}\setlength\itemsep{6pt}
  \item \textbf{Pocos tipos:} \texttt{INTEGER} para números enteros, \texttt{REAL} para decimales y \texttt{TEXT} para todo lo demás, incluso las fechas.
  \item \textbf{No hay} \texttt{CREATE DATABASE}: la base de datos \textbf{es el archivo}.
  \item Lo demás (\texttt{PRIMARY KEY}, \texttt{NOT NULL}, \texttt{UNIQUE}) se escribe \textbf{igual}.
\end{itemize}
\vspace{0.15cm}
\small\color{utpGris} \textbf{¿Por qué?} SQLite quiere ser pequeño y simple: con tres tipos le alcanza.
\end{frame}
% ---------------------------------------------------------------------
\begin{frame}{\texttt{FOREIGN KEY} en SQLite: hay que encenderla}
\begin{columns}[c,onlytextwidth]
\column{.52\textwidth}
\begin{sqlcode}[\textwidth]
\scriptsize
\nuevo{\kw{PRAGMA} foreign\_keys = ON;}\\[5pt]
\kw{CREATE TABLE} pedido (\\
\ind id INTEGER \kw{PRIMARY KEY},\\
\ind fecha TEXT,\\
\ind cliente\_id INTEGER,\\
\ind \kw{FOREIGN KEY} (cliente\_id)\\
\ind\ind \kw{REFERENCES} cliente(id)\\
);
\end{sqlcode}
\column{.44\textwidth}
\begin{itemize}\setlength\itemsep{8pt}
  \item Se escribe \textbf{igual} que en MySQL.
  \item Pero viene \textbf{apagada}: hay que encenderla con el \texttt{PRAGMA}.
\end{itemize}
\end{columns}
\vspace{0.15cm}
\begin{center}
\small
\renewcommand{\arraystretch}{1.2}
\begin{tabular}{>{\raggedright\arraybackslash}p{6.0cm}>{\raggedright\arraybackslash}p{6.0cm}}
\toprule
\bfseries Un pedido del cliente 99, que no existe & \bfseries Qué pasa \\
\midrule
Sin el \texttt{PRAGMA} & \mal{Se guarda}, y apunta a nadie \\
Con el \texttt{PRAGMA} & SQLite lo \textbf{rechaza} \\
\bottomrule
\end{tabular}
\end{center}
\vspace{0.15cm}
\analogia{Es como una \textbf{alarma} que viene instalada pero apagada: si no la enciendes, no avisa.}
\end{frame}
% ---------------------------------------------------------------------
\begin{frame}{\texttt{ALTER TABLE} en SQLite: solo lo básico}
\begin{columns}[T,onlytextwidth]
\column{.52\textwidth}
\centering\etiqueta{Sí se puede}
\begin{sqlcode}[\textwidth]
\scriptsize
\com{agregar una columna}\\
\kw{ALTER TABLE} cliente\\
\ind \nuevo{\kw{ADD COLUMN}} telefono TEXT;\\[5pt]
\com{cambiarle el nombre a una columna}\\
\kw{ALTER TABLE} cliente\\
\ind \nuevo{\kw{RENAME COLUMN}} correo \kw{TO} email;
\end{sqlcode}
\column{.44\textwidth}
\centering\etiqueta{No se puede}\\[10pt]
\raggedright
\begin{itemize}\setlength\itemsep{8pt}
  \item \textbf{Cambiar el tipo} de una columna.
  \item \textbf{Agregar una clave foránea} a una tabla que ya existe.
\end{itemize}
\end{columns}
\vspace{0.3cm}
\ojo{En SQLite, las claves foráneas se escriben \textbf{al crear la tabla}. Después ya no se pueden agregar.}
\vspace{0.1cm}
\small\color{utpGris} \textbf{¿Por qué?} Para seguir siendo pequeño, SQLite solo trae los cambios más sencillos.
\end{frame}
% ---------------------------------------------------------------------
\begin{frame}{¿Y si necesito justo lo que no se puede?}
Se \textbf{rehace la tabla}, como quien pasa un cuaderno en limpio:
\vspace{0.4cm}
\begin{center}
\begin{tikzpicture}
  \tikzset{paso/.style={cajaverde,font=\small,text width=2.45cm,minimum height=1.35cm}}
  \node[paso] (a) {\textbf{1.} Crear una tabla nueva, corregida};
  \node[paso,right=6mm of a] (b) {\textbf{2.} Copiar las filas};
  \node[paso,right=6mm of b] (c) {\textbf{3.} Borrar la tabla vieja};
  \node[paso,right=6mm of c] (d) {\textbf{4.} Ponerle el nombre de siempre};
  \draw[flecha] (a)--(b); \draw[flecha] (b)--(c); \draw[flecha] (c)--(d);
\end{tikzpicture}
\end{center}
\vspace{0.5cm}
\concepto{Lo más fácil mientras aprendes}{\small Corrige el \texttt{CREATE TABLE} en tu archivo \texttt{.sql} y \textbf{crea la base de nuevo}. Para eso sirve el archivo que veremos al final.}
\end{frame}
% ---------------------------------------------------------------------
\begin{frame}{\texttt{DROP TABLE} en SQLite}
\begin{sqlcode}[.60\textwidth]
\kw{DROP TABLE} pedido;
\end{sqlcode}
\vspace{-0.1cm}
Se escribe \textbf{igual} que en MySQL: borra la tabla \textbf{con todos sus datos}.
\vspace{0.3cm}
\begin{center}
\renewcommand{\arraystretch}{1.35}
\begin{tabular}{>{\raggedright\arraybackslash}p{3.8cm}>{\ttfamily\small}p{4.1cm}>{\raggedright\arraybackslash}p{4.3cm}}
\toprule
\bfseries Quiero\ldots & \normalfont\etiqueta{MySQL} & \etiqueta{SQLite} \\
\midrule
Borrar una tabla & DROP TABLE pedido; & \texttt{\small DROP TABLE pedido;} \\
Vaciar una tabla & TRUNCATE TABLE pedido; & \texttt{\small DELETE FROM pedido;} \\
Borrar toda la base & DROP DATABASE tienda; & \textbf{borrar el archivo} \\
\bottomrule
\end{tabular}
\end{center}
\vspace{0.25cm}
\ojo{Borra primero la tabla que \textbf{apunta} (\tbl{pedido}) y después la \textbf{apuntada} (\tbl{cliente}).}
\end{frame}
% ---------------------------------------------------------------------
\begin{frame}{\ins{Práctica 4} De MySQL a SQLite}
Este SQL salió de dbdiagram para MySQL. Escribe cómo queda \textbf{cada línea} en SQLite.
\vspace{0.3cm}
\begin{center}
\renewcommand{\arraystretch}{1.1}
\begin{tabular}{>{\ttfamily\footnotesize}p{6.4cm}>{\raggedright\arraybackslash}p{5.8cm}}
\toprule
\normalfont\bfseries En MySQL & \bfseries En SQLite \\
\midrule
CREATE DATABASE colegio; & \bl[5.5cm] \\[10pt]
id INT PRIMARY KEY AUTO\_INCREMENT, & \bl[5.5cm] \\[10pt]
nombre VARCHAR(80) NOT NULL, & \bl[5.5cm] \\[10pt]
fecha DATE NOT NULL, & \bl[5.5cm] \\[10pt]
TRUNCATE TABLE matricula; & \bl[5.5cm] \\
\bottomrule
\end{tabular}
\end{center}
\vspace{0.3cm}
\small \textbf{Pista:} una de las líneas \textbf{no existe} en SQLite. ¿Y qué línea falta al comienzo para encender las claves foráneas?
\end{frame}
% =======================================================================
\seccion{Archivos .sql}
% ---------------------------------------------------------------------
\begin{frame}{Un archivo \texttt{.sql}: instrucciones guardadas}
\begin{columns}[c,onlytextwidth]
\column{.54\textwidth}
\etiqueta{tienda.sql}
\vspace{-0.15cm}
\begin{sqlcode}[\textwidth]
\scriptsize
\com{Tienda: tablas y datos de prueba}\\[4pt]
\kw{CREATE TABLE} cliente (\\
\ind id INTEGER \kw{PRIMARY KEY},\\
\ind nombre TEXT \kw{NOT NULL}\\
);\\[4pt]
\kw{INSERT INTO} cliente (nombre) \kw{VALUES} ('Ana');\\
\kw{INSERT INTO} cliente (nombre) \kw{VALUES} ('Luis');
\end{sqlcode}
\column{.42\textwidth}
\begin{itemize}\setlength\itemsep{8pt}
  \item Es un archivo de \textbf{texto} con instrucciones SQL.
  \item El gestor las ejecuta \textbf{de arriba hacia abajo}.
  \item Se abre con \textbf{cualquier editor}.
\end{itemize}
\end{columns}
\vspace{0.3cm}
Es lo mismo que has escrito hasta hoy, pero \textbf{guardado} para no perderlo.
\end{frame}
% ---------------------------------------------------------------------
\begin{frame}{¿Para qué sirve un archivo \texttt{.sql}?}
\begin{columns}[T,onlytextwidth]
\column{.31\textwidth}
\centering\etiqueta{Repetir}\\[6pt]
\raggedright\small Crea la base completa \textbf{desde cero}, las veces que haga falta.
\column{.31\textwidth}
\centering\etiqueta{Compartir}\\[6pt]
\raggedright\small Lo envías y la otra persona obtiene \textbf{la misma} base.
\column{.31\textwidth}
\centering\etiqueta{Respaldar}\\[6pt]
\raggedright\small Queda una \textbf{copia} de las tablas y de los datos.
\end{columns}
\vspace{0.5cm}
\begin{center}
\begin{tikzpicture}
  \node[cajaverde,text width=3.4cm] (s) {\texttt{tienda.sql}\\[-1pt]{\scriptsize las \textbf{instrucciones}}};
  \node[cajanaranja,text width=3.4cm,right=26mm of s] (d) {\texttt{tienda.db}\\[-1pt]{\scriptsize la \textbf{base de datos}}};
  \draw[flecha] (s)--(d) node[midway,above,font=\scriptsize] {se ejecuta};
\end{tikzpicture}
\end{center}
\vspace{0.3cm}
\analogia{El \texttt{.sql} es la \textbf{receta}; la base de datos es \textbf{el plato}. Con la receta lo vuelves a preparar cuando quieras.}
\end{frame}
% ---------------------------------------------------------------------
\begin{frame}[plain]
  \begin{tikzpicture}[remember picture,overlay]
    \node[anchor=north west,font=\bfseries\fontsize{24}{28}\selectfont]
      at ([xshift=12mm,yshift=-9mm]current page.north west) {Preguntas};
    \node[anchor=north east,inner sep=0pt]
      at ([xshift=-10mm,yshift=-8mm]current page.north east)
      {\includegraphics[width=43mm]{utp-logo.png}};
    \node[anchor=center,inner sep=0pt]
      at ([yshift=-5mm]current page.center)
      {\includegraphics[width=78mm]{preguntas-ilustracion.png}};
  \end{tikzpicture}
\end{frame}
% ---------------------------------------------------------------------
\begin{frame}{Cierre: lo de hoy en una tabla}
\begin{center}
\renewcommand{\arraystretch}{1.5}
\begin{tabular}{>{\raggedright\bfseries}p{3.4cm}>{\raggedright\arraybackslash}p{9.0cm}}
\toprule
& \bfseries En una frase \\
\midrule
Vista & Una consulta guardada con nombre. No guarda datos. \\
Diagrama como código & Un texto que un programa convierte en dibujo. \\
DBML & El plano de la base de datos: ayuda a SQL, no lo reemplaza. \\
DDL en SQLite & Lo mismo de siempre, más simple: pocos tipos y un archivo. \\
Archivo \texttt{.sql} & La receta para crear la base de datos otra vez. \\
\bottomrule
\end{tabular}
\end{center}
\vspace{0.4cm}
\small El camino completo: \textbf{piensas} en DBML, lo \textbf{ves} como diagrama, lo \textbf{exportas} a SQL, lo \textbf{guardas} en un \texttt{.sql} y lo \textbf{ejecutas}.
\end{frame}
% ---------------------------------------------------------------------
\begin{frame}{Referencias}
\footnotesize
Coronel, C., Morris, S., \& Rob, P. (2011). \textit{Bases de datos: diseño, implementación y administración} (9.ª ed.). Cengage Learning.

\vspace{0.3cm}
Elmasri, R., \& Navathe, S. (2007). \textit{Fundamentos de sistemas de bases de datos} (5.ª ed.). Addison Wesley.

\vspace{0.3cm}
Silberschatz, A., Korth, H. F., \& Sudarshan, S. (2020). \textit{Database system concepts} (7th ed.). McGraw-Hill.

\vspace{0.3cm}
Holistics. (s.\,f.). \textit{DBML: Database Markup Language}. \texttt{https://dbml.dbdiagram.io}

\vspace{0.3cm}
SQLite. (s.\,f.). \textit{SQL as understood by SQLite}. \texttt{https://www.sqlite.org/lang.html}
\end{frame}
% ---------------------------------------------------------------------
% Cierre UTP
{
  \setbeamertemplate{background canvas}{%
    \includegraphics[width=\paperwidth,height=\paperheight]{gracias-final.pdf}%
  }
  \setbeamertemplate{footline}{}
  \begin{frame}[plain]
  \end{frame}
}

\end{document}
